\chapter{Procedimentos metodológicos}
Descreva o enquadramento metodológico da pesquisa~\cite{gil2022}, o objeto de estudo, os instrumentos de coleta de dados e as etapas de análise. A Figura~\ref{fig:roteiro} apresenta um exemplo editável de roteiro metodológico desenhado diretamente em TikZ (fluxo sequencial composto por coleta de dados, etapas 1 a 4 e análise dos resultados).

\begin{figure}[htbp]
  \centering
  \caption{Roteiro de trabalho}\label{fig:roteiro}
  \begin{tikzpicture}[
    node distance=0.55cm,
    etapa/.style={
      rectangle,
      rounded corners=2pt,
      draw=blue!55!black,
      line width=0.8pt,
      minimum width=8cm,
      text width=7.2cm,
      minimum height=0.72cm,
      align=center,
      font=\small
    },
    seta/.style={
      -{Triangle[length=2.2mm,width=3.2mm]},
      draw=blue!45!gray!60,
      line width=1.6pt
    }
  ]
    \node[etapa] (e0) {Coleta de dados};
    \node[etapa, below=of e0] (e1) {Etapa 1};
    \node[etapa, below=of e1] (e2) {Etapa 2};
    \node[etapa, below=of e2] (e3) {Etapa 3};
    \node[etapa, below=of e3] (e4) {Etapa 4};
    \node[etapa, below=of e4] (e5) {Análise dos resultados};

    \draw[seta] (e0) -- (e1);
    \draw[seta] (e1) -- (e2);
    \draw[seta] (e2) -- (e3);
    \draw[seta] (e3) -- (e4);
    \draw[seta] (e4) -- (e5);
  \end{tikzpicture}
  \fonte{Elaboração própria}
\end{figure}

\noindent\textbf{Descrição do roteiro:} o processo segue em sequência, da coleta de dados às etapas 1, 2, 3 e 4, e termina na análise dos resultados.
